\chapter{\ph{Título del apéndice}}
\label{app:apendice_a}

\ph{Los apéndices recogen material de referencia que rompería la lectura del
cuerpo principal pero que conviene tener disponible: especificaciones largas,
ejemplos completos de artefactos generados, manuales de uso, datasets,
representaciones gráficas o textuales del modelo del dominio, etc.}

\ph{Sugerencias de contenido típico:}

\begin{itemize}
  \item \ph{Especificación formal completa (gramática, esquema, modelo).}
  \item \ph{Ejemplo extenso de un artefacto que produce el sistema.}
  \item \ph{Manual de instalación / despliegue / uso.}
  \item \ph{Tablas de referencia muy extensas.}
  \item \ph{Figuras grandes que requieren orientación apaisada.}
\end{itemize}

\ph{Si el TFG necesita varios apéndices, duplica este archivo
(\texttt{B\_apendice.tex}, \texttt{C\_apendice.tex}\dots) y añade
los correspondientes \texttt{\textbackslash include} en \texttt{main.tex}
dentro del bloque \texttt{\textbackslash appendix}.}

\ph{Para incluir bloques de código o ficheros largos, usa el entorno
\texttt{verbatim} (ya configurado con numeración de líneas):}

\begin{verbatim}
[Reemplaza este bloque por el contenido literal que quieras incluir:
 fragmento de código, fichero de configuración, ejemplo completo, etc.]
\end{verbatim}
